% Einheit 2 von 3 – Potenzgesetze (bank/reelle-zahlen/e2.jsonl, zusammenbau v0.1)
%% TODO Zweigzeile Teil 1: was hier gelernt wird (Satz in Schülersprache aus dem Einheitstitel) – Einheit 2
\einheitenkopf[e2]{Einheit 2 von 3 · Potenzgesetze}
\zweigzeile{ab Kl. 9, je nach Buch bis Kl. 10 · keine P10-Aufgabe}
\setcounter{aufgabe}{13}

%% TODO Titel: Ich-kann-Satz fehlt in der Bank (Platzhalter: Kettenname) – Nr. 14
%% TODO Anweisung: ein Satz über der ersten Teilaufgabe fehlt in der Bank – Nr. 14
\begin{aufgabe}{Potenzgesetze}
%% TODO \rechenplatz{n}: Schrittzahl des Grundfalls fehlt in der Bank (nur bei fester Schreibform, 2.3 b)
\begin{teile}
\teil $7^3 \cdot 7^5$ – welches Gesetz passt? \\ \kreuz{gleiche Basis} \\ \kreuz{gleicher Exponent} \\ \kreuz{keins: ausrechnen}
\teil $3^2 \cdot 3^3$ – als Malkette, wie viele Faktoren, als Potenz? \leerfeld
\teil $5^4 \cdot 5^3$ – als eine Potenz und als Zahl? \leerfeld
\teil $7^8 : 7^3$ – als eine Potenz und als Zahl? \leerfeld
\teil $(2^3)^2$ – als eine Potenz und als Zahl? \leerfeld
\teil $2^4 \cdot 5^4$ – als eine Potenz und als Zahl? \leerfeld
\teil $3^2 \cdot 2^3$ – mit Gesetz oder ausrechnen? Ergebnis? \leerfeld
\teil $3^2 \cdot 3^2 \cdot 2$ – Ergebnis? \leerfeld
\teil $x^5 \cdot x^2$ – vereinfacht? \leerfeld
\teil $2x^3 \cdot 5x^4$ – vereinfacht? \leerfeld
\teil $3^4 + 3^4 + 3^4$ – als eine Potenz und als Zahl? \leerfeld
\teil $(3xy)^2$ – vereinfacht? \leerfeld
\teil $\frac{6x^5 \cdot 2x^3}{4x^2}$ – vereinfacht? \leerfeld
\end{teile}
\end{aufgabe}

%% TODO Titel: Ich-kann-Satz fehlt in der Bank (Platzhalter: Kettenname) – Nr. 15
%% TODO Anweisung: ein Satz über der ersten Teilaufgabe fehlt in der Bank – Nr. 15
\begin{aufgabe}{Potenzen mit negativem Exponenten in Brüche und zurück (Blatt 0 aus potenzen-wurzeln.md; LISUM-PH)}
\begin{teile}
\teil $5^{-2}$ – als Bruch? \leerfeld
\end{teile}
\end{aufgabe}

%% TODO Titel: Ich-kann-Satz fehlt in der Bank (Platzhalter: Kettenname) – Nr. 16
%% TODO Anweisung: ein Satz über der ersten Teilaufgabe fehlt in der Bank – Nr. 16
\begin{aufgabe}{Potenzgesetze – Fehler finden, Begründen, Anwendung}
\begin{teile}
\teil Finde den Fehler und rechne richtig: \rechnung{3^2 \cdot 3^5 &= 9^7}
\teil Warum ist $3^2 \cdot 3^4$ nicht $9^6$? Begründe mit der Malkette.
\teil Ein Speicherbaustein fasst $2^{30}$ Byte. Ein Laufwerk enthält $2^8$ solcher Bausteine. Wie viele Byte fasst das Laufwerk? Gib eine Zweierpotenz an.
\end{teile}
\end{aufgabe}
